\begin{frame}{Why Stage 1 Is a Rare-Event Problem}
  \begin{columns}[T,onlytextwidth]
    \column{0.48\textwidth}
      \centering
      {\bfseries\large Typical balanced\\classification\par}
      \vspace{0.55em}
      \begin{minipage}{0.92\linewidth}
        \begin{itemize}
          \item Positive and negative classes are closer in size
          \item Accuracy and ROC-style metrics are often informative
          \item The goal is usually overall classification correctness
        \end{itemize}
      \end{minipage}

    \column{0.48\textwidth}
      \centering
      {\color{orange!90!black}\bfseries\large Our Stage 1 task\par}
      \vspace{1.2em}
      \begin{minipage}{0.92\linewidth}
        \begin{itemize}
          \item Positive events are very rare (\(<5\%\))
          \item We care about ranking the few important positives near the top
          \item \textbf{PR-AUC} is more useful than plain accuracy here
        \end{itemize}
      \end{minipage}
  \end{columns}

  \vspace{0.5em}
  \small
  Intuition: PR-AUC focuses on whether the model's high-risk predictions are truly positive and whether rare positive cases are captured early.
\end{frame}
